\subsection{Ảnh hưởng của cấu trúc encoding}

Tại $N=64$, Sequential có 47{,}752 clause và trung vị
18{,}143 giây; Pairwise có 428{,}864 clause nhưng trung vị chỉ
0{,}322 giây. Cả hai hoàn thành 5/5 lần. Đây là một đối chứng rõ
ràng chống lại cách chọn encoding chỉ bằng số clause. Ít clause hơn
không bảo đảm tìm kiếm nhanh hơn vì công thức với biến phụ tạo cấu trúc
lan truyền và lựa chọn nhánh khác. Ở $N=128$, Product có 192{,}890
clause với trung vị 0{,}719 giây, còn Pairwise có 3{,}462{,}784
clause và trung vị 3{,}063 giây. Như vậy quan hệ giữa kích thước CNF
và thời gian thay đổi theo cả encoding lẫn kích thước.

Kết quả Product cũng không đơn điệu: toàn bộ năm run ở $N=32,64$
đều timeout với ngưỡng 60 giây, nhưng cả năm run ở $N=128$ hoàn
thành và cả ba run ở $N=1000$ hoàn thành với ngưỡng 120 giây.
Giới hạn thời gian ở phase Large dài hơn, nên không thể đặt hai
phase này thành một đường cong cùng điều kiện. Ngay trong cùng phase
Standard, việc Product timeout ở $N=64$ rồi thành công ở $N=128$
cho thấy kích thước bàn cờ không phải biến giải thích duy nhất; heuristic
và cấu trúc CNF có thể tương tác phức tạp. Dữ liệu không đủ để nhận diện
một cơ chế nhân quả duy nhất.

Ở $N=1000$, Product tạo 1{,}298{,}508 biến và 11{,}942{,}164
clause. Thời gian dựng CNF trung vị là 7{,}736 giây; phần gọi solver,
bao gồm nạp clause, có trung vị 20{,}391 giây; trung vị tổng thời gian
là 28{,}127 giây. Hai thành phần không nên cộng như những trung vị
độc lập để suy ra chính xác trung vị tổng, nhưng chúng cho thấy cả
chi phí dựng lẫn chi phí solver đều đáng kể ở kích thước này.

\subsection{Chi tiết 13 run Gurobi sau khi cấp Academic WLS}

Bảng~\ref{tab:full-gurobi} giữ từng lượt Gurobi bổ sung, thay vì chỉ
trình bày một trung vị. Tất cả 13 run đều có nghiệm hợp lệ. Tại
$N=128$, bốn run dưới 0{,}58 giây và một run 0{,}879 giây;
trung vị 0{,}194 giây thể hiện nhóm thời gian thấp nhưng không che
giấu biến thiên giữa các lần đo. Tại $N=256$, khoảng thời gian là
1{,}539--3{,}654 giây. Sự biến thiên này là một lý do không dựa
vào một run đơn lẻ để kết luận ưu thế của solver.

\begin{table}[htbp]
\centering
\caption{Thời gian từng run Gurobi Academic WLS tại các kích thước bổ sung.
Các run đều có trạng thái SAT và nghiệm hợp lệ.}
\label{tab:full-gurobi}
\input{\ReportPrefix generated/full_gurobi_runs.tex}
\end{table}

Mô hình Gurobi tại $N=64,128,256$ lần lượt có $4{,}096$,
$16{,}384$ và $65{,}536$ biến nhị phân, cùng $382$, $766$ và
$1{,}534$ ràng buộc. Những số ràng buộc khớp $6N-2$ từ mô hình
đã trình bày. So sánh thời gian Gurobi với SAT/CP trong cùng bảng
chỉ mang tính tham chiếu: các run Academic WLS được tạo ở phiên
Colab khác, dù cùng pipeline và giới hạn thời gian tương ứng.

\subsection{Hai cách dùng CPLEX trong thực nghiệm}

Bảng~\ref{tab:full-cplex} tách CPLEX MIP và CPLEX CP Optimizer.
Đây là hai solver khác mô hình và khác hệ cấp phép, không phải hai
tên cho cùng một run. CP Optimizer giải thành công 3/3 lần ở
$N=8,16,32$ với mô hình $N$ biến nguyên và ba
\texttt{AllDifferent}. CPLEX MIP giải thành công 3/3 lần ở
$N=8,16$, nhưng cả ba run $N=32$ gặp \texttt{LICENSE\_LIMIT}
trước khi có thống kê thời gian giải. Gurobi Academic WLS không
mở khóa license CPLEX MIP.

\begin{table}[htbp]
\centering
\caption{Kết quả hai solver IBM trên Colab. Dấu -- là không có trung vị
thời gian của run hoàn thành.}
\label{tab:full-cplex}
\input{\ReportPrefix generated/full_cplex_table.tex}
\end{table}

Hai unit test CPLEX có thể bị skip khi chạy trên máy local nếu thiếu
thư viện \texttt{cplex}/\texttt{docplex} hoặc executable
\texttt{cpoptimizer}. Việc skip đó phản ánh dependency trên máy local,
không phủ nhận 18 run CPLEX trong dữ liệu Colab: 15 SAT hợp lệ và ba
\texttt{LICENSE\_LIMIT}. Khi trình bày kết quả, cần giữ ranh giới
giữa trạng thái kiểm thử local và kết quả thực nghiệm đã lưu.
